\section{Описание}
Требуется преобразовать натуральное число $n$ в единицу с помощью операций <<вычесть единицу>>, <<разделить на два>> и <<разделить на три>> (деление только нацело). Стоимость каждой операции равна текущему значению числа, а стоимость преобразования~--- сумме стоимостей всех операций. Нужно найти преобразование минимальной стоимости и вывести его стоимость и последовательность операций.

Задача решается методом динамического программирования. Жадный выбор здесь не подходит: например, для $n = 10$ деление на два даёт число 5, из которого выгоднее вычесть единицу, чем делить дальше, а локально лучший шаг не всегда ведёт к глобальному оптимуму. При этом выполняется свойство оптимальной подструктуры: если оптимальное преобразование числа $i$ начинается с некоторой операции, то остаток пути является оптимальным преобразованием полученного числа.

Пусть $dp[i]$~--- минимальная стоимость преобразования числа $i$ в единицу. Тогда
\begin{equation*}
dp[1] = 0, \qquad
dp[i] = i + \min
\begin{cases}
dp[i-1], \\
dp[i/2], & \text{если } i \bmod 2 = 0, \\
dp[i/3], & \text{если } i \bmod 3 = 0.
\end{cases}
\end{equation*}
Слагаемое $i$ появляется потому, что первая операция применяется к числу $i$ и стоит ровно $i$. Значения вычисляются в порядке возрастания $i$, так как все $dp[j]$, от которых зависит $dp[i]$, имеют $j < i$ и уже посчитаны.

Для восстановления ответа вместе со стоимостью в каждой ячейке хранится операция, на которой достигается минимум. После вычисления таблицы начинаем с $i = n$ и пока $i > 1$ выводим сохранённую операцию и переходим к $i-1$, $i/2$ или $i/3$ соответственно. Операции выводятся в порядке их применения, так что дополнительного разворота не требуется.

Сложность по времени~--- $O(n)$, так как каждое значение $dp[i]$ вычисляется за $O(1)$. Сложность по памяти~--- тоже $O(n)$: при $n \le 10^7$ таблица из элементов по 8 байт занимает около 80~МБ.

\pagebreak

\section{Исходный код}
Операции описаны перечислением $operaion$ с тремя значениями: $minus\_1$, $div\_2$ и $div\_3$. Для ячейки таблицы создана структура $elem$, где хранятся стоимость $price$ и операция $op$, на которой достигнут минимум.

Таблица $dp$ заполняется для $i$ от 2 до $n$. Сначала берётся вариант <<вычесть единицу>>, затем, если $i$ делится на 2 или на 3 и соответствующая стоимость строго меньше текущей, ячейка обновляется вариантом с делением. В конце к стоимости прибавляется $i$. Базовые элементы $dp[0]$ и $dp[1]$ имеют стоимость 0; $dp[0]$ не используется и нужен лишь для того, чтобы индексы совпадали со значениями $n$.

\begin{lstlisting}[language=C++]
#include <iostream>
#include <vector>

enum operaion {
    minus_1,
    div_2,
    div_3
};

struct elem {
    int price = 0;
    operaion op;

    elem() {}
    elem(int p, operaion o) : price(p), op(o) {}
};

int main(void) {
    int n;
    std::cin >> n;

    std::vector<elem> dp {elem{}, elem{}};
    dp.reserve(n+1);

    for (int i = 2; i <= n; ++i) {
        dp.emplace_back(dp[i-1].price, minus_1);
        if (i % 2 == 0 && dp[i / 2].price < dp[i].price) {
            dp[i] = {dp[i / 2].price, div_2};
        }
        if (i % 3 == 0 && dp[i / 3].price < dp[i].price) {
            dp[i] = {dp[i / 3].price, div_3};
        }
        dp[i].price += i;
    }

    std::cout << dp[n].price << '\n';

    for (int i = n; i > 1;) {
        if (dp[i].op == minus_1) {
            std::cout << "-1 ";
            --i;
        } else if (dp[i].op == div_2) {
            std::cout << "/2 ";
            i /= 2;
        } else if (dp[i].op == div_3) {
            std::cout << "/3 ";
            i /= 3;
        }
    }
}
\end{lstlisting}

\section{Консоль}
\begin{alltt}
laba-7> cmake -S . -B build
laba-7> cmake --build build
laba-7> ..build/Debug/main.exe
10
21
/2 -1 /2 -1
\end{alltt}

Проверка для $n = 10$: операции $10 \to 5 \to 4 \to 2 \to 1$ стоят $10 + 5 + 4 + 2 = 21$.
\pagebreak